% GENERATED from paper/draft/01_introduction.md by build.py — do not edit
\section{Introduction}

\subsection{The delegation question}

Organisations that deploy AI agents must decide which decisions an agent may take without human review. Supervisors now
ask them to make this decision explicitly and to revisit it. The Financial Stability Board's 2026 sound practices
suggest that firms ``may also progressively increase an agent's ability to interact with the external environment as
they gain confidence it can do so safely'', with ``triggers for transitioning between different forms of human oversight
based on performance metrics or user requests'' (\citealp{fsb2026sound}). The UK Financial Conduct Authority (FCA) calls for
``continuous, conditional assurance'' and for frameworks covering ``delegated authority'' (\citealp{rusu2026fca}). Neither
says what evidence should ground the confidence on which autonomy is extended.

\subsection{Why track records are not enough}

The default answer is a track record: success rates, human approvals, correction rates, model confidence, or a judge's
reading of the agent's trace. A track record answers whether the agent's actions turned out correct. A delegator also
needs to know whether they were correct \emph{for the right reasons and within scope}, because only then can past
success be expected to transfer to the next case. This is the premise of the paper. The pilot tested one consequence of
it, that a rule built on evidence tests, that is, on interventions on the evidence the agent consults, delegates more
safely than a track record, and §6.1 reports that it could not confirm this. Outcome success masks the basis of a
decision, much as human correction masks an agent's raw competence (\citealp{azevedo2026minimal}).

The gap is not hypothetical. In CXBench, 80.1\% of runs that passed the outcome check failed process checks in regulated
customer-experience workflows (\citealp{yu2026cxbench}). Procedure-aware evaluation finds that between 27\% and 78\% of
reported τ-bench successes, depending on model and domain, are ``corrupt successes'' (\citealp{cao2026pae}). In paired tasks,
an agent's skill at acting correctly and its skill at correctly abstaining are nearly independent (phi coefficient φ ≈
−0.10; \citealp{liu2026agentabstain}). Self-reports do not close the gap: verbalised confidence barely separates correct from
incorrect answers (area under the ROC curve, AUROC, 0.47--0.66; \citealp{shen2026provable}), and stated reasoning can be
unfaithful to what drives a model's answer (\citealp{turpin2023unfaithful}).

There is also a formal reason. Relational properties of an agent's behaviour, such as what its action depends on, are
2-safety hyperproperties that no monitor of a single trace can decide; deciding them requires comparing executions
(\citealp{pandey2026local}). Calibrating on final outcomes cannot control the share of correct outcomes reached on an
unsupported basis (\citealp{opoku2026toolchain}). The philosophy of meaningful human control reached the same conclusion from
another direction: a system must ``demonstrably and verifiably be responsive to the human moral reasons relevant in the
circumstances'', and this holds ``no matter how efficient {[}it{]} may be in performing specific tasks''
(\citealp{santonidesio2018mhc}).

\subsection{A motivating example}

We met this gap in our own work. Auditable Design, an open-source pipeline that turns user reviews into design briefs,
promises that ``every finding has a typed chain back to informing user reviews''. In its data, the typed chain ends at
the level of a review cluster (27 reviews for the showcased cluster). The link from an individual finding to specific
reviews exists only as quote indices inside the free-text reasoning. The record is rich and readable, and it is not
wrong. But from the record alone a reviewer cannot tell which reviews each finding depends on. Observable is not the
same as accountable.

\subsection{Question, finding and contributions}

We ask which evidence should decide that an agent may take a class of decisions without review: its track record, an LLM
judge's reading of its trace, or its reaction to changes in the evidence. For the last we define
\emph{decision-structured agency} (DSA): the extent to which an agent configuration, within a class of decisions,
regularly recognises and reacts to the evidence that ought to settle those decisions. It also covers escalating when
that evidence is missing or contradictory or when an action would exceed the agent's authority. The notion descends from
reasons-responsiveness (\citealp{fischer1998responsibility}) and from the tracking condition of meaningful human control
(\citealp{santonidesio2018mhc}).

Our pre-registered pilot does not show that delegating on DSA beats a track record. The registered test favoured DSA,
but the advantage rested on cases where our own policy admits more than one reading, and it disappeared when those cases
were excluded or the answer key was reviewed. Under every answer key, an LLM judge passed most of the decisions that the
evidence tests marked as resting on the wrong basis. The main lesson concerns method: in an evaluation of this kind the
answer key is part of the instrument. It feeds both the test and the criterion against which the test is judged, and two
separately written oracles (programs that compute the correct action from the generator's hidden facts) that agree with
each other can still share the same contestable reading. \citet{cao2026loopholes} shows that policy ambiguity can
masquerade as agent error in benchmark scores; we show that it can also decide which delegation rule appears to work.

We make three contributions.

\textbf{C1 --- A construct and its measurement.} We define DSA per pair (agent configuration, decision class) and
measure it by composed interventions on the external evidence the agent consults, with control arms, scoring responses
against the agent-declared grounds (§3). Which response each intervention requires is fixed by oracles that encode one
reading of the policy; the pilot shows how much rests on that reading.

\textbf{C2 --- A delegation rule.} We give a rule that grants, narrows and revokes autonomy per decision class beneath a
fixed, human-set ceiling, with certificates bound to the agent's configuration (§4). The statistical machinery is
standard; the novelty is the loss it controls. The pilot exercises only the granting step; narrowing, revocation and
transfer between classes are specified but not yet tested.

\textbf{C3 --- A pre-registered pilot with a full audit trail.} On a synthetic counterparty-verification workflow with
14 model configurations, we compare delegation on DSA with delegation on a track record, the registered primary
comparison, and with rules built on an LLM judge, at matched autonomy and under evidence shift (§5, §6). We publish the
registration, every amendment, the log of policy ambiguities, the author's review of the oracles' answers and the code,
so that each step from data to conclusion can be checked.

\subsection{Scope}

We study ordinary, non-adversarial shifts in evidence. When sources can be corrupted adversarially, no procedure that
sees the same evidence can avoid handing off ambiguous cases to a human (\citealp{pandey2026certified}); for that setting we
propose combining our rule with a source-quorum gate. Our guarantees hold for generated episodes and relative to a
pre-registered family of interventions. What the rule provides is reliability allocation beneath a fixed ceiling, not a
security guarantee. The case study encodes publicly available anti-money-laundering rules in a fully synthetic
environment. It uses no real entities and no institutional data, and makes no claim of regulatory compliance.

\begin{center}\rule{0.5\linewidth}{0.5pt}\end{center}
